\documentclass[11pt,a4paper,openany]{book}

% =============================================================================
% Professional textbook foundation
% =============================================================================

\usepackage[T1]{fontenc}
\usepackage[utf8]{inputenc}
\usepackage{lmodern}
\usepackage{microtype}
\usepackage[a4paper,margin=27mm,headheight=15pt]{geometry}
\usepackage{setspace}
\setstretch{1.05}
\setlength{\parindent}{1.15em}
\setlength{\parskip}{0pt}
\raggedbottom

\usepackage{amsmath,amssymb,mathtools}
\usepackage{booktabs,tabularx,array,longtable}
\usepackage{graphicx}
\usepackage{listings}

\usepackage{tikz}
\usepackage{pgfplots}
\usepackage{caption}
\usepackage{subcaption}
\usepackage{float}
\usepackage{wrapfig}
\usepackage{adjustbox}
\usetikzlibrary{
  arrows.meta,
  angles,
  quotes,
  calc,
  positioning,
  intersections,
  decorations.pathmorphing,
  decorations.markings,
  backgrounds,
  fit,
  shapes.geometric,
  patterns,
  3d
}
\pgfplotsset{compat=1.18}

\newcolumntype{Y}{>{\raggedright\arraybackslash}X}

\usepackage{xcolor}
\definecolor{chapterblue}{RGB}{35,75,120}
\definecolor{softblue}{RGB}{240,246,252}
\definecolor{softgreen}{RGB}{241,249,244}
\definecolor{softgold}{RGB}{252,248,236}
\definecolor{softred}{RGB}{252,242,242}
\definecolor{softgray}{RGB}{246,247,249}
\definecolor{rulegray}{RGB}{185,190,198}

\lstdefinestyle{prbookpython}{
  language=Python,
  basicstyle=\ttfamily\small,
  keywordstyle=\bfseries,
  commentstyle=\itshape,
  stringstyle=\ttfamily,
  numbers=left,
  numberstyle=\scriptsize,
  numbersep=8pt,
  frame=single,
  rulecolor=\color{rulegray},
  backgroundcolor=\color{softgray},
  breaklines=true,
  showstringspaces=false,
  columns=fullflexible,
  keepspaces=true
}
\lstdefinestyle{prbookwolfram}{
  basicstyle=\ttfamily\small,
  numbers=left,
  numberstyle=\scriptsize,
  numbersep=8pt,
  frame=single,
  rulecolor=\color{rulegray},
  backgroundcolor=\color{softgray},
  breaklines=true,
  showstringspaces=false,
  columns=fullflexible,
  keepspaces=true
}

\usepackage[most]{tcolorbox}
\tcbset{
  enhanced,
  breakable,
  boxrule=0.6pt,
  arc=1.4mm,
  left=2mm,right=2mm,top=1.5mm,bottom=1.5mm,
  before skip=7pt,after skip=7pt
}

\newtcolorbox{keyidea}[1][Key Idea]{
  title=#1,colback=softblue,colframe=chapterblue,fonttitle=\bfseries
}
\newtcolorbox{definitionbox}[1][Definition]{
  title=#1,colback=softgreen,colframe=green!45!black,fonttitle=\bfseries
}
\newtcolorbox{physicalbox}[1][Physical Interpretation]{
  title=#1,colback=softgold,colframe=orange!55!black,fonttitle=\bfseries
}
\newtcolorbox{mistakebox}[1][Common Mistake]{
  title=#1,colback=softred,colframe=red!55!black,fonttitle=\bfseries
}
\newtcolorbox{historybox}[1][Historical Note]{
  title=#1,colback=softgray,colframe=rulegray,fonttitle=\bfseries
}
\newtcolorbox{examplebox}[1][Worked Example]{
  title=#1,colback=white,colframe=chapterblue,fonttitle=\bfseries
}
\newtcolorbox{solutionbox}[1][Solution]{
  title=#1,colback=softgray,colframe=rulegray,fonttitle=\bfseries
}
\newtcolorbox{roadmapbox}[1][Chapter Roadmap]{
  title=#1,colback=softgray,colframe=chapterblue,fonttitle=\bfseries
}

\usepackage{enumitem}
\setlist[itemize]{leftmargin=1.8em,topsep=4pt,itemsep=2pt,parsep=0pt,partopsep=0pt}
\setlist[enumerate]{leftmargin=2em,topsep=4pt,itemsep=3pt,parsep=0pt}

\usepackage{titlesec}
\titleformat{\chapter}[display]
  {\normalfont\bfseries\color{chapterblue}}
  {\filleft\Large\chaptername\ \thechapter}
  {0.7ex}
  {\titlerule\vspace{1.1ex}\Huge}
  [\vspace{0.7ex}\titlerule]
\titlespacing*{\chapter}{0pt}{-16pt}{20pt}

\titleformat{\section}
  {\Large\bfseries\color{chapterblue}}
  {\thesection}{0.7em}{}
\titlespacing*{\section}{0pt}{2.1ex plus .5ex minus .2ex}{0.7ex}

\titleformat{\subsection}
  {\large\bfseries}
  {\thesubsection}{0.7em}{}
\titlespacing*{\subsection}{0pt}{1.7ex plus .4ex minus .2ex}{0.5ex}

\usepackage{fancyhdr}
\pagestyle{fancy}
\fancyhf{}
\fancyhead[LE]{\small\itshape Theory of Quantum Mechanics}
\fancyhead[RO]{\small\itshape\nouppercase{\rightmark}}
\fancyfoot[C]{\thepage}
\renewcommand{\headrulewidth}{0.4pt}
\renewcommand{\footrulewidth}{0pt}
\renewcommand{\chaptermark}[1]{\markboth{\thechapter.\ #1}{}}
\renewcommand{\sectionmark}[1]{\markright{\thesection.\ #1}}

\usepackage{tocloft}
\setcounter{tocdepth}{2}
\setcounter{secnumdepth}{2}
\renewcommand{\contentsname}{Table of Contents}
\renewcommand{\cftchapfont}{\bfseries}
\renewcommand{\cftchappagefont}{\bfseries}
\setlength{\cftbeforechapskip}{7pt}
\setlength{\cftbeforesecskip}{1pt}
\renewcommand{\cftdotsep}{1.5}

\usepackage[hidelinks,hypertexnames=false]{hyperref}

\newcommand{\chapterquote}[2]{%
  \begin{center}
    \begin{minipage}{0.86\textwidth}
      \centering\itshape\small #1
      \if\relax\detokenize{#2}\relax
      \else
        \par\vspace{0.45em}\raggedleft\normalfont\footnotesize--- #2
      \fi
    \end{minipage}
  \end{center}
  \vspace{0.7em}
}

\newcommand{\learningobjectives}{%
  \section*{Learning Objectives}
  \addcontentsline{toc}{section}{Learning Objectives}
  \noindent After completing this chapter, the reader should be able to:
}

\newcommand{\chaptersummary}{%
  \section*{Chapter Summary}
  \addcontentsline{toc}{section}{Chapter Summary}
}

\newcommand{\nextchapter}{%
  \section*{Looking Ahead}
  \addcontentsline{toc}{section}{Looking Ahead}
}




%%%%%%%%%%%%%%%%%%%%%%%%%%%%%%%%%%%%%%%%%%%%%%%%%%%%%%%%%%%%%%%%%%%%%%%%%%%%%%%%
% PR-BOOK-12 : Editorial hierarchy and short-concept environments
%%%%%%%%%%%%%%%%%%%%%%%%%%%%%%%%%%%%%%%%%%%%%%%%%%%%%%%%%%%%%%%%%%%%%%%%%%%%%%%%
% Numbered hierarchy: Part -> Chapter -> Section -> Subsection.
% Shorter concepts use an unnumbered editorial heading or a semantic box.
\newcommand{\booktopic}[1]{%
  \par\addvspace{1.05\baselineskip}%
  \noindent{\large\bfseries #1}\par\nobreak
  \vspace{0.28\baselineskip}%
}

\newtcolorbox{propertybox}[1][Property]{
  enhanced,breakable,colback=blue!2,colframe=blue!45!black,
  fonttitle=\bfseries,title=#1,boxrule=0.45pt,arc=1.2mm,
  left=1.5mm,right=1.5mm,top=1mm,bottom=1mm
}
\newtcolorbox{rulebox}[1][Rule]{
  enhanced,breakable,colback=cyan!3,colframe=cyan!45!black,
  fonttitle=\bfseries,title=#1,boxrule=0.45pt,arc=1.2mm,
  left=1.5mm,right=1.5mm,top=1mm,bottom=1mm
}
\newtcolorbox{remarkbox}[1][Remark]{
  enhanced,breakable,colback=black!2,colframe=black!35,
  fonttitle=\bfseries,title=#1,boxrule=0.4pt,arc=1.2mm,
  left=1.5mm,right=1.5mm,top=1mm,bottom=1mm
}

\newenvironment{property}[1][Property]{\begin{propertybox}[#1]}{\end{propertybox}}
\newenvironment{bookrule}[1][Rule]{\begin{rulebox}[#1]}{\end{rulebox}}
\newenvironment{bookremark}[1][Remark]{\begin{remarkbox}[#1]}{\end{remarkbox}}


%%%%%%%%%%%%%%%%%%%%%%%%%%%%%%%%%%%%%%%%%%%%%%%%%%%%%%%%%%%%%%%%%%%%%%%%%%%%%%%%
% PR-BOOK-02 : Professional textbook environments
%%%%%%%%%%%%%%%%%%%%%%%%%%%%%%%%%%%%%%%%%%%%%%%%%%%%%%%%%%%%%%%%%%%%%%%%%%%%%%%%

\newcounter{definition}[chapter]
\renewcommand{\thedefinition}{\thechapter.\arabic{definition}}

\newcounter{example}[chapter]
\renewcommand{\theexample}{\thechapter.\arabic{example}}

\newcounter{remark}[chapter]
\renewcommand{\theremark}{\thechapter.\arabic{remark}}

\newenvironment{definition}{
\refstepcounter{definition}
\begin{definitionbox}[Definition \thedefinition]
}{
\end{definitionbox}
}

\newenvironment{workedexample}{
\refstepcounter{example}
\begin{examplebox}[Example \theexample]
}{
\end{examplebox}
}

\newenvironment{solution}{
\begin{solutionbox}[Solution]
}{
\end{solutionbox}
}

\newenvironment{physical}{
\begin{physicalbox}[Physical Interpretation]
}{
\end{physicalbox}
}

\newenvironment{warningbox}{
\begin{mistakebox}[Common Mistake]
}{
\end{mistakebox}
}

\newenvironment{historicalnote}{
\begin{historybox}[Historical Note]
}{
\end{historybox}
}

\newcommand{\ImportantEquation}[1]{
\begin{keyidea}[Key Equation]
\[
#1
\]
\end{keyidea}
}



%%%%%%%%%%%%%%%%%%%%%%%%%%%%%%%%%%%%%%%%%%%%%%%%%%%%%%%%%%%%%%%%%%%%%%%%%%%%%%%%
% PR-BOOK-03 : Complete educational framework
%%%%%%%%%%%%%%%%%%%%%%%%%%%%%%%%%%%%%%%%%%%%%%%%%%%%%%%%%%%%%%%%%%%%%%%%%%%%%%%%

% Additional counters ---------------------------------------------------------
\newcounter{keyequation}[chapter]
\renewcommand{\thekeyequation}{\thechapter.\arabic{keyequation}}

\newcounter{exercise}[chapter]
\renewcommand{\theexercise}{\thechapter.\arabic{exercise}}

\newcounter{historicalnote}[chapter]
\renewcommand{\thehistoricalnote}{\thechapter.\arabic{historicalnote}}

% Additional colours ----------------------------------------------------------
\definecolor{summaryblue}{RGB}{232,241,250}
\definecolor{notationgray}{RGB}{248,249,251}
\definecolor{exercisegreen}{RGB}{239,248,241}
\definecolor{outcomegold}{RGB}{253,248,229}
\definecolor{lookaheadviolet}{RGB}{247,242,251}

% Educational boxes -----------------------------------------------------------
\newtcolorbox{keyequationbox}[1]{
  title={Key Equation \thekeyequation: #1},
  colback=softblue,
  colframe=chapterblue,
  fonttitle=\bfseries,
  borderline west={2.2pt}{0pt}{chapterblue}
}

\newtcolorbox{summarybox}[1][Chapter Summary]{
  title=#1,
  colback=summaryblue,
  colframe=chapterblue,
  fonttitle=\bfseries\large
}

\newtcolorbox{formulasheetbox}[1][Quick Formula Sheet]{
  title=#1,
  colback=white,
  colframe=chapterblue,
  fonttitle=\bfseries\large,
  boxrule=0.8pt
}

\newtcolorbox{notationbox}[1][Notation and Units]{
  title=#1,
  colback=notationgray,
  colframe=rulegray,
  fonttitle=\bfseries\large
}

\newtcolorbox{outcomesbox}[1][Learning Outcomes Check]{
  title=#1,
  colback=outcomegold,
  colframe=orange!60!black,
  fonttitle=\bfseries\large
}

\newtcolorbox{lookingaheadbox}[1][Looking Ahead]{
  title=#1,
  colback=lookaheadviolet,
  colframe=violet!55!black,
  fonttitle=\bfseries\large
}

\newtcolorbox{exercisebox}[2][]{
  title={Exercise \theexercise\if\relax\detokenize{#2}\relax\else: #2\fi},
  colback=exercisegreen,
  colframe=green!45!black,
  fonttitle=\bfseries,
  #1
}

\newtcolorbox{dashboardbox}[1][Chapter Dashboard]{
  title=#1,
  colback=softgray,
  colframe=rulegray,
  fonttitle=\bfseries\large
}

% Semantic environments -------------------------------------------------------
\newenvironment{keyequation}[1]{%
  \refstepcounter{keyequation}%
  \begin{keyequationbox}{#1}%
}{%
  \end{keyequationbox}%
}

\newenvironment{chapterexercise}[1][]{%
  \refstepcounter{exercise}%
  \begin{exercisebox}{#1}%
}{%
  \end{exercisebox}%
}

\newenvironment{chapterhistorical}[1][]{%
  \refstepcounter{historicalnote}%
  \begin{historybox}[Historical Note \thehistoricalnote%
  \if\relax\detokenize{#1}\relax\else: #1\fi]%
}{%
  \end{historybox}%
}

% Structured worked-example helpers ------------------------------------------
\newcommand{\problempart}{\par\medskip\noindent\textbf{Problem.}\ }
\newcommand{\givenpart}{\par\medskip\noindent\textbf{Given.}\ }
\newcommand{\findpart}{\par\medskip\noindent\textbf{Find.}\ }
\newcommand{\solutionpart}{\par\medskip\noindent\textbf{Solution.}\ }
\newcommand{\verificationpart}{\par\medskip\noindent\textbf{Verification.}\ }
\newcommand{\meaningpart}{\par\medskip\noindent\textbf{Physical meaning.}\ }

% Difficulty tags -------------------------------------------------------------
\newcommand{\difficulty}[1]{%
  \hfill\textbf{Difficulty: }%
  \ifcase#1\or$\star\circ\circ\circ\circ$%
  \or$\star\star\circ\circ\circ$%
  \or$\star\star\star\circ\circ$%
  \or$\star\star\star\star\circ$%
  \or$\star\star\star\star\star$\fi
}

% Consistent chapter-ending headings ------------------------------------------
\newcommand{\chapterformulaheading}{%
  \section*{Quick Formula Sheet}
  \addcontentsline{toc}{section}{Quick Formula Sheet}
}
\newcommand{\chapternotationheading}{%
  \section*{Notation and Units}
  \addcontentsline{toc}{section}{Notation and Units}
}
\newcommand{\chapteroutcomesheading}{%
  \section*{Learning Outcomes Check}
  \addcontentsline{toc}{section}{Learning Outcomes Check}
}
\newcommand{\chapterexercisesheading}{%
  \section*{Exercises}
  \addcontentsline{toc}{section}{Exercises}
}
\newcommand{\chapterlookaheadheading}{%
  \section*{Looking Ahead}
  \addcontentsline{toc}{section}{Looking Ahead}
}

% Compact reusable tables -----------------------------------------------------
\newenvironment{notationtable}{%
  \begin{tabular}{>{\centering\arraybackslash}p{0.16\textwidth}
                  >{\raggedright\arraybackslash}p{0.48\textwidth}
                  >{\raggedright\arraybackslash}p{0.23\textwidth}}
  \toprule
  \textbf{Symbol} & \textbf{Meaning} & \textbf{SI unit}\\
  \midrule
}{%
  \bottomrule
  \end{tabular}
}

\newenvironment{formulatable}{%
  \begin{tabular}{>{\raggedright\arraybackslash}p{0.29\textwidth}
                  >{\raggedright\arraybackslash}p{0.61\textwidth}}
  \toprule
  \textbf{Quantity or operation} & \textbf{Formula}\\
  \midrule
}{%
  \bottomrule
  \end{tabular}
}



%%%%%%%%%%%%%%%%%%%%%%%%%%%%%%%%%%%%%%%%%%%%%%%%%%%%%%%%%%%%%%%%%%%%%%%%%%%%%%%%
% PR-BOOK-04 : Visual identity and TikZ illustration system
%%%%%%%%%%%%%%%%%%%%%%%%%%%%%%%%%%%%%%%%%%%%%%%%%%%%%%%%%%%%%%%%%%%%%%%%%%%%%%%%

% Visual palette ----------------------------------------------------------------
\definecolor{visualblue}{RGB}{35,75,120}
\definecolor{visualcyan}{RGB}{54,135,170}
\definecolor{visualgold}{RGB}{193,139,46}
\definecolor{visualred}{RGB}{164,63,63}
\definecolor{visualgreen}{RGB}{58,126,86}
\definecolor{visualink}{RGB}{38,45,54}
\definecolor{visualgrid}{RGB}{210,218,226}

% Global TikZ styles ------------------------------------------------------------
\tikzset{
  physics axis/.style={
    -{Latex[length=3mm]},
    line width=0.9pt,
    draw=visualink
  },
  physics vector/.style={
    -{Latex[length=3.2mm]},
    line width=1.5pt,
    draw=visualblue
  },
  physics vector secondary/.style={
    -{Latex[length=3mm]},
    line width=1.25pt,
    draw=visualcyan
  },
  force vector/.style={
    -{Latex[length=3.2mm]},
    line width=1.45pt,
    draw=visualred
  },
  field vector/.style={
    -{Latex[length=2.8mm]},
    line width=1.1pt,
    draw=visualgreen
  },
  construction line/.style={
    dashed,
    line width=0.75pt,
    draw=visualgrid
  },
  trajectory/.style={
    line width=1.5pt,
    draw=visualblue
  },
  visual label/.style={
    font=\small,
    text=visualink,
    fill=white,
    inner sep=1.5pt
  },
  concept node/.style={
    rounded corners=2pt,
    draw=visualblue,
    fill=softblue,
    line width=0.8pt,
    align=center,
    inner sep=5pt,
    font=\small
  }
}

% Figure caption system ---------------------------------------------------------
\captionsetup{
  font=small,
  labelfont={bf,color=visualblue},
  textfont={color=visualink},
  justification=raggedright,
  singlelinecheck=false,
  skip=6pt
}

\InputIfFileExists{figures/common/chapter_figure_styles.tex}{}{}

\newcommand{\figsource}[1]{%
  \par\smallskip
  {\footnotesize\color{gray}\textit{Source/Construction:} #1}
}

\newenvironment{bookfigure}[2][tbp]{%
  \def\bookfigurecaption{#2}%
  \gdef\bookfigurelabel{}%
  \let\bookfigureoriginallabel\label
  \renewcommand{\label}[1]{\gdef\bookfigurelabel{##1}}%
  \begin{figure}[#1]
  \centering
  \begin{adjustbox}{max width=\textwidth}
}{%
  \end{adjustbox}
  \caption{\bookfigurecaption}%
  \if\relax\detokenize\expandafter{\bookfigurelabel}\relax\else
    \bookfigureoriginallabel{\bookfigurelabel}%
  \fi
  \end{figure}
}

% Compact visual callout --------------------------------------------------------
\newtcolorbox{visualguidebox}[1][How to read the figure]{
  title=#1,
  colback=white,
  colframe=visualcyan,
  fonttitle=\bfseries,
  borderline west={2pt}{0pt}{visualcyan}
}

% Chapter-opening artwork -------------------------------------------------------
\newcommand{\chapteropeningart}[2]{%
\begin{center}
\begin{tikzpicture}[x=1cm,y=1cm]
  \fill[softblue] (0,0) rectangle (12.6,2.65);
  \draw[visualblue,line width=1pt] (0,0) rectangle (12.6,2.65);
  \draw[visualgrid,line width=0.5pt] (0.45,0.35) grid[xstep=0.55,ystep=0.55] (12.1,2.3);
  \node[anchor=west,font=\bfseries\Large,text=visualblue] at (0.7,1.78) {#1};
  \node[anchor=west,align=left,text width=7.2cm,font=\small,text=visualink]
    at (0.72,0.88) {#2};
  \draw[physics axis] (9.1,0.62) -- (11.85,0.62) node[right] {$x$};
  \draw[physics axis] (9.1,0.62) -- (9.1,2.18) node[above] {$y$};
  \draw[physics vector] (9.1,0.62) -- (11.08,1.73)
    node[above right,visual label] {$\mathbf q$};
  \fill[visualgold] (9.1,0.62) circle (1.7pt);
\end{tikzpicture}
\end{center}
\vspace{0.4cm}
}

% Future-ready visual templates -------------------------------------------------
\newcommand{\wavefunctionplot}{%
\begin{tikzpicture}
  \begin{axis}[
    width=10.5cm,
    height=5.2cm,
    axis lines=middle,
    xlabel={$x$},
    ylabel={$\psi(x)$},
    ticks=none,
    domain=-3.3:3.3,
    samples=180,
    ymin=-0.75,
    ymax=1.15
  ]
    \addplot[visualblue,very thick] {exp(-x^2/2)};
    \addplot[visualcyan,thick,dashed] {0.72*x*exp(-x^2/2)};
  \end{axis}
\end{tikzpicture}
}

\newcommand{\quantumhalltemplate}{%
\begin{tikzpicture}[x=1cm,y=1cm]
  \fill[softblue] (0,0) rectangle (9.8,4.3);
  \draw[visualblue,line width=0.9pt] (0,0) rectangle (9.8,4.3);
  \foreach \x in {0.7,1.7,...,9.2}
    \foreach \y in {0.6,1.6,...,3.8}
      \draw[visualcyan!55,line width=0.5pt] (\x,\y) circle (0.16);
  \draw[force vector] (0.55,3.75) -- (9.15,3.75)
    node[midway,above,visual label] {upper chiral edge state};
  \draw[force vector] (9.15,0.55) -- (0.55,0.55)
    node[midway,below,visual label] {lower chiral edge state};
  \node[font=\small] at (4.9,2.15) {2DEG bulk: Landau-level states};
  \draw[physics vector] (8.85,1.5) -- (8.85,2.85)
    node[above] {$\mathbf B$};
\end{tikzpicture}
}

\title{\Huge\bfseries Theory of Quantum Mechanics\\[0.55em]
\Large From Mathematical Foundations to the Quantum Hall Effect}
\author{}
\date{PR-BOOK v1.8 - Chapters 1--2 Publication Foundation\\\today}

% =============================================================================
% PR-BOOK-08.0 publishing services
% =============================================================================
\usepackage{amsthm}
\usepackage[noautomatic]{imakeidx}
\makeindex[title=Subject Index,columns=2]

\usepackage[acronym,toc,nonumberlist]{glossaries}
\makeglossaries

% Bibliography is optional until citations are populated.
\usepackage[numbers,sort&compress]{natbib}

\newtheorem{theorem}{Theorem}[chapter]
\newtheorem{lemma}[theorem]{Lemma}
\newtheorem{proposition}[theorem]{Proposition}
\newtheorem{corollary}[theorem]{Corollary}
\theoremstyle{definition}
\newtheorem{postulate}[theorem]{Postulate}
\newtheorem{principle}[theorem]{Principle}
\theoremstyle{remark}
\newtheorem{remarktheorem}[theorem]{Remark}

\newcommand{\term}[1]{\textbf{#1}\index{#1}}
\newcommand{\notation}[2]{\ensuremath{#1} & #2\\}

%%%%%%%%%%%%%%%%%%%%%%%%%%%%%%%%%%%%%%%%%%%%%%%%%%%%%%%%%%%%%%%%%%%%%%%%%%%%%%%%
% PR-BOOK-13 : Repository cleanup and publication typography
%%%%%%%%%%%%%%%%%%%%%%%%%%%%%%%%%%%%%%%%%%%%%%%%%%%%%%%%%%%%%%%%%%%%%%%%%%%%%%%%
% PR-BOOK-13: Canonical typography, hierarchy, and page furniture.
% Loaded after the historical preamble so these settings are authoritative.

% Paragraph rhythm
\setstretch{1.055}
\setlength{\parindent}{1.15em}
\setlength{\parskip}{0pt}
\setlength{\emergencystretch}{2em}
\widowpenalty=10000
\clubpenalty=10000
\displaywidowpenalty=10000
\brokenpenalty=7000
\predisplaypenalty=6000
\postdisplaypenalty=0

% Numbering and contents depth: Part -> Chapter -> Section -> Subsection.
\setcounter{secnumdepth}{2}
\setcounter{tocdepth}{2}

% Part openings
\titleformat{\part}[display]
  {\centering\normalfont\bfseries\color{chapterblue}}
  {\Large\MakeUppercase{\partname}\ \thepart}
  {1.2ex}
  {\titlerule[0.8pt]\vspace{1.4ex}\Huge}
\titlespacing*{\part}{0pt}{3.5cm}{2.5cm}

% Chapter and heading rhythm
\titleformat{\chapter}[display]
  {\normalfont\bfseries\color{chapterblue}}
  {\filleft\Large\chaptername\ \thechapter}
  {0.7ex}
  {\titlerule[0.8pt]\vspace{1.1ex}\Huge}
\titlespacing*{\chapter}{0pt}{-14pt}{22pt}

\titleformat{\section}
  {\Large\bfseries\color{chapterblue}}
  {\thesection}{0.72em}{}
\titlespacing*{\section}{0pt}{2.35ex plus .55ex minus .25ex}{0.8ex}

\titleformat{\subsection}
  {\large\bfseries\color{black!88}}
  {\thesubsection}{0.72em}{}
\titlespacing*{\subsection}{0pt}{1.9ex plus .4ex minus .2ex}{0.55ex}

% Unnumbered editorial topic used below subsection level.
\renewcommand{\booktopic}[1]{%
  \par\addvspace{1.08\baselineskip}%
  \noindent{\large\bfseries\color{black!88}#1}\par\nobreak
  \vspace{0.3\baselineskip}%
}

% Running heads: book title on verso, current chapter/section on recto.
\fancypagestyle{mainmatterstyle}{%
  \fancyhf{}
  \fancyhead[LE]{\small\scshape\nouppercase{\leftmark}}
  \fancyhead[RO]{\small\itshape\nouppercase{\rightmark}}
  \fancyfoot[LE,RO]{\thepage}
  \renewcommand{\headrulewidth}{0.35pt}
  \renewcommand{\footrulewidth}{0pt}
}
\pagestyle{mainmatterstyle}
\fancypagestyle{plain}{%
  \fancyhf{}
  \fancyfoot[C]{\thepage}
  \renewcommand{\headrulewidth}{0pt}
  \renewcommand{\footrulewidth}{0pt}
}
\renewcommand{\chaptermark}[1]{\markboth{\chaptername\ \thechapter.\ #1}{}}
\renewcommand{\sectionmark}[1]{\markright{\thesection.\ #1}}

% TOC hierarchy and spacing
\renewcommand{\cftpartfont}{\large\bfseries\color{chapterblue}}
\renewcommand{\cftpartpagefont}{\large\bfseries\color{chapterblue}}
\setlength{\cftbeforepartskip}{14pt}
\setlength{\cftbeforechapskip}{6.5pt}
\setlength{\cftbeforesecskip}{0.7pt}
\renewcommand{\cftchapfont}{\bfseries}
\renewcommand{\cftchappagefont}{\bfseries}
\renewcommand{\cftdotsep}{1.6}

% PR-BOOK-13: Canonical semantic environment aliases.
% These aliases establish one vocabulary while preserving legacy environments.

\newenvironment{bookdefinition}[1][Definition]
  {\begin{definitionbox}[#1]}
  {\end{definitionbox}}

\newenvironment{bookproperty}[1][Property]
  {\begin{propertybox}[#1]}
  {\end{propertybox}}

\newenvironment{bookprinciple}[1][Principle]
  {\begin{keyidea}[#1]}
  {\end{keyidea}}

\newenvironment{bookexample}[1][Example]
  {\begin{examplebox}[#1]}
  {\end{examplebox}}

\newenvironment{bookinsight}[1][Physical Insight]
  {\begin{physicalbox}[#1]}
  {\end{physicalbox}}

\newenvironment{bookhistory}[1][Historical Note]
  {\begin{historybox}[#1]}
  {\end{historybox}}

\newenvironment{bookwarning}[1][Common Mistake]
  {\begin{mistakebox}[#1]}
  {\end{mistakebox}}

\newenvironment{booksummary}[1][Summary]
  {\begin{summarybox}[#1]}
  {\end{summarybox}}

% PR-BOOK-13: Cross-reference and PDF publication services.
\usepackage{bookmark}
\usepackage[nameinlink,noabbrev,capitalise]{cleveref}
\usepackage{xurl}

\crefname{chapter}{Chapter}{Chapters}
\crefname{section}{Section}{Sections}
\crefname{subsection}{Section}{Sections}
\crefname{figure}{Figure}{Figures}
\crefname{table}{Table}{Tables}
\crefname{equation}{Equation}{Equations}
\crefname{theorem}{Theorem}{Theorems}
\crefname{lemma}{Lemma}{Lemmas}
\crefname{proposition}{Proposition}{Propositions}
\crefname{corollary}{Corollary}{Corollaries}
\crefname{definition}{Definition}{Definitions}
\crefname{example}{Example}{Examples}
\crefname{exercise}{Exercise}{Exercises}

\hypersetup{
  pdftitle={Theory of Quantum Mechanics},
  pdfsubject={Graduate textbook in mathematical and quantum physics},
  pdfkeywords={quantum mechanics, mathematical physics, classical mechanics, vector calculus},
  pdfcreator={LaTeX},
  pdfproducer={PR-BOOK-13 publication system},
  bookmarksopen=true,
  bookmarksnumbered=true,
  bookmarksdepth=2,
  pdfdisplaydoctitle=true,
  colorlinks=true,
  linkcolor=chapterblue,
  citecolor=green!35!black,
  urlcolor=chapterblue
}

